\documentclass[../cuaderno6.tex]{subfiles}

\begin{document}

\subsection{Disco}

En el caso del disco, el campo coercitivo, la magnetización remanente y máxima varían de forma coherente entre sí con el ángulo de medición. Se puede ver en las figuras \ref{fig:H_c_segun_H_a_Disco_Fa0_27_08}, \ref{fig:M_r_segun_H_a_Disco_Fa0_27_08} y \ref{fig:M_max_segun_H_a_Disco_Fa0_27_08} como la variación respeta las simetrías del disco frente a ser rotado más de 90 grados, obteniendo resultados similares para cuando la rotación sólo equivale a una inversión en el signo del campo aplicado.

\begin{figure}[H]
    \centering
    \includegraphics[width=0.5\textwidth]{../img/Fa0/Disco/H_c_segun_H_a_Disco_Fa0_27_08.png}
    \caption{El campo coercitivo $H_c$ en función del logaritmo del campo aplicado $H_{max}$ para cada ángulo.}
    \label{fig:H_c_segun_H_a_Disco_Fa0_27_08}
\end{figure}
\begin{figure}[h]
    \centering
    \includegraphics[width=0.5\textwidth]{../img/Fa0/Disco/M_r_segun_H_a_Disco_Fa0_27_08.png}
    \caption{Magnetización remanente $M_r$ en función del campo aplicado $H_{max}$.}
    \label{fig:M_r_segun_H_a_Disco_Fa0_27_08}
\end{figure}
\begin{figure}[h]
    \centering
    \includegraphics[width=0.5\textwidth]{../img/Fa0/Disco/M_max_segun_H_a_Disco_Fa0_27_08.png}
    \caption{Magnetización máxima $M_{max}$ en función del campo aplicado $H_{max}$.}
    \label{fig:M_max_segun_H_a_Disco_Fa0_27_08}
\end{figure}

El factor de desmagnetización $N_d$ también sigue con la tendencia general de variar apropiadamente para cada rotación ejecutada, visible en la figura \ref{fig:N_d_angulo_major_Disco} corrobora la bajísima susceptibilidad magnética del disco en ciertas orientaciones, como con su cara más extensa de frente al campo magnético. Es decir, con su normal en paralelo al campo. En este caso, se registra correctamente un pico marcado en el factor de desmagnetización correspondiente a la superficie reducida para que cierren las líneas del campo magnético.

\begin{figure}[H]
    \centering
    \includegraphics[width=0.5\textwidth]{../img/Fa0/Disco/N_d_angulo_major_Disco.png}
    \caption{Factor de desmagnetización $N_d$ en función del ángulo.}
    \label{fig:N_d_angulo_major_Disco}
\end{figure}

En el caso del ajuste para el campo coercitivo, visible en la figura \ref{fig:H0_c_angulo_minors_Disco}, es notorio observar como este resulta menor para el caso de 90 grados que sus adyacentes.

\begin{figure}[H]
    \centering
    \includegraphics[width=0.5\textwidth]{../img/Fa0/Disco/H0_c_angulo_minors_Disco.png}
    \caption{Gráfico de la pendiente del ajuste $H_c^0$ en función del ángulo.}
    \label{fig:H0_c_angulo_minors_Disco}
\end{figure}

\end{document}